% ==============================================================================
% Chapter 6: Conclusion, Research Trends, and Open Gaps
% Owner: Member 2 (TV2)
% ==============================================================================

\chapter{Kết luận, xu hướng và khoảng trống nghiên cứu}
\label{chap:conclusion}

\section{Tổng kết các phát hiện chính của báo cáo}
\label{sec:conc_summary}

Bài khảo sát \textit{Ask in Any Modality} đặt Multimodal RAG trong
một khung gồm các lựa chọn về truy xuất, kết hợp thông tin, tăng cường
ngữ cảnh, tạo sinh và huấn luyện. Khung này giúp xác định công đoạn mà
mỗi nghiên cứu cải thiện và điều kiện để chuyển kết quả của nghiên
cứu vào một hệ thống khác \cite{abootorabi2025ask}. Ba thông điệp ở
slide kết luận của nhóm tập trung vào đường đi của chứng cứ từ kho
dữ liệu đến câu trả lời.

\paragraph{Ảnh trang giữ chứng cứ về bố cục.}
ColPali và VisRAG cho thấy trang PDF có thể được lập chỉ mục trực
tiếp dưới dạng ảnh. Trên ViDoRe, ColPali cải thiện nDCG@5 trung bình
từ 66,1 của BGE-M3 + OCR lên 81,3. Các mức tăng lớn ở hình khoa học
và infographic cho thấy có những trường hợp việc giữ bố cục ảnh
làm thay đổi đáng kể kết quả tìm trang. Ở tầng sinh, ví dụ LinkedIn
trong bài VisRAG cho thấy cả hai pipeline có thể chọn đúng cùng
một trang nhưng chỉ cấu hình đọc ảnh gắn đúng tỷ lệ với vùng địa
lý. Cần đánh giá tìm trang và đọc chứng cứ bằng các phép đo riêng
\cite{faysse2025colpali,yu2025visrag}.

Khả năng đọc bố cục có chi phí. ColPali lưu khoảng 1030 vector
128 chiều cho một trang và so khớp từng token truy vấn. VisRAG-Ret
lưu một vector 2304 chiều, thuận tiện cho chỉ mục vector. Điểm
MRR@10 của VisRAG-Ret nhỉnh hơn trung bình trong phép đối chiếu
sáu tập, còn ColPali có lợi thế trên DocVQA và SlideVQA trong
cùng thiết lập. Một hệ thống dùng nhiều vector có thể đầu tư vào
độ phân giải và khả năng kiểm tra vùng khớp. Một hệ thống dùng
vector đơn có thể dùng phần tài nguyên tiết kiệm để duy trì kho
lớn hơn hoặc cấp thêm ứng viên cho bộ đọc. Việc chọn cấu hình cần
gắn với truy vấn thực tế, ngân sách bộ nhớ và yêu cầu độ trễ
\cite{faysse2025colpali,yu2025visrag}.

\paragraph{Định tuyến theo phương thức và độ chi tiết.}
Một kho có văn bản, bảng, ảnh và video đặt ra câu hỏi trước khi
xếp hạng. Hệ thống phải xác định loại dữ liệu có thể trả lời
và mức chi tiết thích hợp. Ví dụ, câu hỏi về một sự kiện ngắn cần
ứng viên ở cấp clip, còn câu hỏi về bố cục của cả một báo cáo có
thể cần tài liệu đầy đủ. UniversalRAG dùng bộ định tuyến chọn kho
theo phương thức và độ chi tiết. Kết quả gần Oracle trong thiết
lập được trình bày ở slide 15 cho thấy tiềm năng của việc chọn
kho trước khi so khớp. Điểm Oracle chỉ là mốc tham chiếu trong
thiết lập đó. Ở miền mới, sai số định tuyến có thể loại chứng cứ
trước khi bộ truy xuất có cơ hội xếp hạng \cite{yeo2026universalrag}.

\paragraph{Đối soát khi chứng cứ không nhất quán.}
Khi bộ sinh có câu trả lời từ tri thức tham số và đồng thời đọc
chứng cứ bên ngoài, hai nguồn có thể dẫn đến các đáp án khác nhau.
Ảnh và văn bản trong tập truy xuất cũng có thể chỉ đến các thực
thể khác nhau. CoRe-MMRAG coi những trường hợp này là hai dạng
xung đột PRKI và VTKI, rồi đối soát qua nhiều bước. Ví dụ máy bay
trong slide đặt ra cả việc nhận đúng thực thể từ ảnh lẫn việc
xác nhận mốc thời gian từ nguồn ngoài. Giá trị của cơ chế đối
soát phụ thuộc vào độ bao phủ và chất lượng của chứng cứ được
truy xuất. Nếu cả tập trang đều bỏ sót chứng cứ cần thiết, nên
ghi nhận thiếu căn cứ để trả lời \cite{tian2025coremmrag}.

Khi ghép các ý trên thành một pipeline, có thể theo dõi lỗi theo
thứ tự. Tầng định tuyến chọn kho, tầng truy xuất chọn chứng cứ
trong kho, tầng đọc và đối soát dùng chứng cứ để tạo đáp án.
Một câu trả lời có thể sai do bất kỳ bước nào, và lỗi ở bước
trước làm giảm khả năng sửa lỗi ở bước sau. Vì vậy, báo cáo độ
chính xác đầu cuối cần đi cùng dấu vết của kho đã chọn, mã trang
đã truy xuất và nguồn dùng để suy luận. Các kết quả nghiên cứu
hiện tại minh họa từng tầng dưới những bộ dữ liệu và chỉ số cụ
thể.

\section{Các xu hướng nghiên cứu chủ đạo}
\label{sec:conc_trends}

\paragraph{Cải thiện VLM làm nền cho truy xuất.}
Slide 17 nêu bước chuyển từ PaliGemma sang Qwen2-VL. Biến thể
ColQwen2 đạt 86,6 nDCG@5 so với 81,3 của ColPali tham chiếu,
tăng 5,3 điểm trong Bảng 7 của nghiên cứu. ColQwen2 dùng chiến
lược huấn luyện tương tự và giới hạn số mảnh ảnh ở 768 để gần
mức bộ nhớ của cấu hình tham chiếu. Đây là bằng chứng về khả
năng thích ứng của cơ chế late interaction khi bộ mã hóa thị
giác--ngôn ngữ mạnh hơn. Số mảnh ảnh và backbone thay đổi cùng
lúc, nên muốn giải thích riêng phần cải thiện do backbone cần
thêm đối chứng \cite{faysse2025colpali}.

Một hệ quả của xu hướng này là nên đánh giá backbone qua kiểu
chứng cứ và điều kiện sử dụng. OCR nội tại của VLM có thể giúp
khớp chữ từ ảnh, nhưng năng lực đó cần đo riêng trên bản quét,
chữ nhỏ và ngôn ngữ ngoài tập huấn luyện. Với biểu đồ, hiểu vị
trí của trục, chú giải và các điểm còn quan trọng hơn tìm được
từ khóa. Hệ thống có thể đạt chất lượng truy xuất tốt mà vẫn
phải thêm bước đọc vùng có độ phân giải cao để trả lời các câu
hỏi về giá trị chính xác.

\paragraph{Duy trì ảnh qua cả truy xuất lẫn tạo sinh.}
VisRAG minh họa xu hướng giữ ảnh làm biểu diễn trung gian từ lập
chỉ mục đến khi chuyển ngữ cảnh cho bộ sinh. Ở đây ảnh là đầu
vào chứng cứ còn câu trả lời có thể ở dạng văn bản. Khái niệm
pipeline thị giác trong slide đề cập đường đi của dữ liệu tài
liệu. Cách này giữ được quan hệ giữa nội dung và bố cục tới khi
bộ sinh đọc trang. Các cơ chế đọc nhiều ảnh mở rộng phép đối
chiếu từ một trang sang nhiều trang \cite{yu2025visrag}.

Một thiết kế ứng dụng có thể giữ mã tài liệu, số trang và ảnh
cùng với biểu diễn dùng để tìm kiếm. Khi phát hiện một kết quả
nhạy với chữ nhỏ, hệ thống có thể chuyển ảnh cắt vùng tương ứng
cho bộ sinh trong khi vẫn bảo toàn thông tin nguồn. Tầng truy
xuất có thể đi từ trang đến vùng, còn tầng sinh cần nối kết quả
với mã trang để người đọc kiểm tra. Những lựa chọn này cần được
đánh giá về độ bao phủ lẫn chi phí xử lý các ảnh bổ sung.

\section{Các khoảng trống nghiên cứu (Research Gaps) còn tồn tại}
\label{sec:conc_gaps}

\subsubsection{Nén chỉ mục đa vector ở quy mô lớn}

Khi một trang có hơn một nghìn vector, chênh lệch dung lượng
nhanh chóng tăng cùng số trang. Nén theo nhóm có thể giảm lặp ở
các vùng nền và giảm chi phí xếp hạng. Nút thắt nằm ở các vùng
nhỏ nhưng chứa đủ chứng cứ, chẳng hạn một ô ghi tỷ lệ trong bảng.
Kết quả token pooling của ColPali cho thấy điểm trung bình giữ
được 97,8\% ở hệ số 3, nhưng các trang dày chữ trong Shift nhạy
hơn với nén \cite{faysse2025colpali}.

Một thí nghiệm tiếp theo có thể chia tập truy vấn theo chứng cứ
cần đọc và thay đổi hệ số gộp theo từng mức. Với mỗi mức, ghi
nDCG, Recall, dung lượng đã lưu, độ trễ truy vấn ở phân vị 95
và tốc độ bổ sung trang mới. Cần ghi thêm trường hợp trang đúng
rơi khỏi danh sách ứng viên sau nén. Tập có nhiều biểu đồ và
tập có nhiều bảng chữ nhỏ nên được báo cáo riêng để xem cùng
một cấu hình phù hợp cho bao nhiêu loại tài liệu.

\subsubsection{Suy luận nhiều trang và nhiều bước}

Một câu hỏi nhiều bước có thể yêu cầu đọc một số liệu, tìm đơn
vị ở trang khác rồi thực hiện phép tính. Nếu chỉ chấm việc tìm
một trang đúng, phép đánh giá có thể bỏ qua việc hệ thống thiếu
phần chứng cứ còn lại. Nếu chỉ chấm kết quả, đáp án đúng do
suy đoán có thể được tính ngang với đáp án kiểm chứng được.
Có thể gán nhãn tập trang cần dùng và kiểm tra cả độ bao phủ
chứng cứ lẫn sự đúng đắn của phép nối thông tin.

Khi tăng số ảnh trong ngữ cảnh, mô hình có thể thấy nhiều chứng
cứ hơn nhưng cũng có thể chọn nhầm trang gần chủ đề. Bảng 3
của VisRAG cho thấy lợi ích trung bình của cấu hình đa ảnh từ
top-1 đến top-3, cùng sự khác biệt theo bộ dữ liệu. Một thí
nghiệm có thể cố định trang đúng và thay đổi số trang nhiễu,
thứ tự trang, độ phân giải và vị trí xuất hiện của chứng cứ.
Kết quả cần ghi cả thời gian sinh, độ chính xác câu trả lời
và tỷ lệ chỉ đúng các trang dùng để suy luận \cite{yu2025visrag}.

\subsubsection{Định vị chứng cứ theo thời gian trong video dài}

Slide 17 đặt temporal grounding như giới hạn của truy xuất thị
giác ở quy mô lớn. Một clip ngắn thuận tiện cho xác định khung
hình chứa sự kiện, nhưng có thể bỏ mất thông tin ngay trước hoặc
sau mốc đó. Cửa sổ dài giữ nhiều bối cảnh và cần nhiều tài nguyên
hơn khi chấm điểm hoặc đưa vào mô hình. UniversalRAG minh họa
việc chọn giữa clip và video ở tầng định tuyến
\cite{yeo2026universalrag}. Sau khi kho đã được chọn, vẫn cần
xác định mốc thời gian thực sự hỗ trợ câu trả lời.

Thử nghiệm trên video dài có thể đánh giá hai tầng. Tầng đầu
truy xuất video hoặc cụm clip. Tầng sau tinh chỉnh vị trí đoạn
liên quan. Cần ghi độ bao phủ video đúng, mức trùng khớp của
đoạn thời gian và chi phí tính toán trên mỗi giờ video.
Với một câu cần hai sự kiện cách xa nhau, điều kiện thành công
phải yêu cầu tìm được cả hai mốc. Cách đánh giá này song song với
bài toán tìm đủ trang trong tài liệu nhiều trang.

\subsubsection{Chuyển miền sang tài liệu tiếng Việt và bản quét}

Kết quả từ ViDoRe và sáu tập của VisRAG tạo cơ sở để chọn
cấu hình thử nghiệm. Khi áp dụng vào tài liệu tiếng Việt,
cần kiểm tra lại khả năng đọc dấu, nhận dạng chữ nhỏ và
liên kết ô bảng với tiêu đề. Chẳng hạn, hồ sơ hành chính
có thể chứa con dấu chồng lên tên đơn vị, còn báo cáo tài
chính có cột đơn vị ghi ở chú thích riêng. Cả hai đều
chứa chữ và bố cục, nhưng cách làm hỏng chứng cứ khác
nhau. Một phép đo theo nhóm bản quét, PDF sinh từ máy và
tài liệu nhiều bảng sẽ cho biết mức độ chuyển miền rõ
hơn so với một điểm trung bình duy nhất.

Tiếng Việt cũng đặt ra yêu cầu về cách viết câu hỏi.
Một mã biểu mẫu có thể xuất hiện cùng cách gọi tắt ở
phần đầu và tên đầy đủ ở phần phụ lục. Người dùng có
thể hỏi bằng tên tắt, cách gõ có dấu hoặc dạng thiếu dấu.
Có thể giữ nguyên trang chứng cứ rồi tạo các biến thể
câu hỏi phù hợp với thói quen tra cứu. Đo vị trí của cùng
một trang qua các biến thể giúp tách lỗi ở biểu diễn
truy vấn khỏi lỗi do độ phân giải và bố cục trang. Đối
với tìm mã, điểm thành công cần xét đúng từng ký tự.
Đối với tìm chủ đề, người đánh giá có thể chấp nhận các
trang khác nếu chúng chứa đủ chứng cứ cho câu trả lời.

Một thiết lập kiểm tra chất lượng ảnh có thể giữ cùng
PDF, trang và truy vấn, rồi so sánh ảnh gốc với phiên
bản giảm độ phân giải hoặc nghiêng nhẹ. Với ColPali, theo
dõi trang đúng và độ ổn định của các vùng đạt điểm
MaxSim cao. Với VisRAG, giữ nguyên danh sách trang
Oracle khi đo khả năng đọc của bộ sinh, sau đó đo lại
trong chế độ truy xuất thực tế. Nếu điểm Oracle giảm khi
thay chất lượng ảnh, lỗi đọc vẫn tồn tại dù đã tìm đúng
trang. Nếu Oracle ổn định nhưng Recall giảm, bước lập
chỉ mục hoặc khớp truy vấn cần được xem xét trước. Cách
bố trí này kiểm tra trực tiếp các giả thuyết về ảnh
trang và lỗi lan truyền ở slide 7, 8 và 16.

Trong một kho thật, lỗi điển hình có thể hiếm nhưng gây
hậu quả lớn, chẳng hạn đọc nhầm một chữ số hoặc đơn vị
trong bảng công nợ. Bộ thử nên ghi đồng thời tỷ lệ tìm
đúng trang, tỷ lệ đọc chính xác giá trị và tỷ lệ dẫn được
chứng cứ. Kết quả cần chỉ ra nhóm nào chịu ảnh hưởng khi
nén chỉ mục hoặc giới hạn số ảnh, rồi đối chiếu mức tiết
kiệm bộ nhớ và thời gian của cấu hình tương ứng. Đây là
một hướng kiểm chứng cho báo cáo, với thí nghiệm gốc của
ColPali và VisRAG làm mốc tham khảo
\cite{faysse2025colpali,yu2025visrag}.

Một câu hỏi chung của ba hướng là hệ thống có thể cung cấp đủ
chứng cứ trong giới hạn thời gian và bộ nhớ hay không. Phép đo
trên một tập truy vấn thống nhất cần ghi rõ kho đã chọn, tập
trang hoặc đoạn đã lấy, chất lượng câu trả lời và thời gian đầu
cuối. Những chỉ số đó cho phép gắn kết quả của từng tầng với
lợi ích quan sát được đối với người đọc.
